% =====================================================================
%  BÁO CÁO TUẦN 2026-W46 — Payment Intelligence Ecosystem
%  Biên dịch: make weekly   (hoặc: xelatex 2026-W46.tex, chạy 2 lần)
%  Hình: ../../slides/figures/ (make slides-figures). Số test chỉ TRÍCH LẠI.
% =====================================================================
\documentclass[11pt,a4paper]{article}

\usepackage{fontspec}
\setmainfont{TeX Gyre Heros}
\setmonofont{DejaVu Sans Mono}[Scale=0.85]
\usepackage[vietnamese]{babel}
\usepackage[a4paper,margin=2.1cm,headheight=14pt]{geometry}
\usepackage{setspace}\setstretch{1.12}
\setlength{\parskip}{0.45em}\setlength{\parindent}{0pt}
\setlength{\emergencystretch}{3em}

\usepackage[table]{xcolor}
\definecolor{navy}{HTML}{1B2A4A}
\definecolor{accent}{HTML}{2A78D6}
\definecolor{accentL}{HTML}{EEF4FC}
\definecolor{warnC}{HTML}{B7791F}
\definecolor{warnL}{HTML}{FDF3E1}
\definecolor{okC}{HTML}{2E8B3E}
\definecolor{okL}{HTML}{E3F5E6}
\definecolor{badC}{HTML}{D03B3B}
\definecolor{ink2}{HTML}{52514E}

\usepackage{graphicx}
\usepackage[most]{tcolorbox}
\usepackage{booktabs,tabularx,array,longtable}
\usepackage{enumitem}\setlist{nosep,leftmargin=1.4em,topsep=0.3em}
\usepackage{titlesec}
\usepackage{fancyhdr}
\usepackage{caption}\captionsetup{font=small,labelfont={bf,color=navy}}
\usepackage{float}
\usepackage{fontawesome5}
\usepackage[hidelinks]{hyperref}
\usepackage{xurl}

\titleformat{\section}{\large\bfseries\color{navy}}{\thesection.}{0.5em}{}
\titlespacing*{\section}{0pt}{1.1em}{0.4em}
\pagestyle{fancy}\fancyhf{}
\fancyhead[L]{\small\color{ink2}Payment Intelligence Ecosystem}
\fancyhead[R]{\small\color{ink2}Báo cáo tuần 2026-W46}
\fancyfoot[C]{\small\thepage}
\renewcommand{\headrulewidth}{0.3pt}

\newtcolorbox{keypoint}[1][Tóm tắt]{enhanced,colback=accentL,colframe=accent,boxrule=0.5pt,arc=2pt,
  left=8pt,right=8pt,top=5pt,bottom=5pt,title={\textbf{#1}},coltitle=white,colbacktitle=accent,
  fonttitle=\small}
\newtcolorbox{warnbox}[1]{enhanced,colback=warnL,colframe=warnC,boxrule=0.5pt,arc=2pt,
  left=8pt,right=8pt,top=5pt,bottom=5pt,title={\textbf{#1}},coltitle=white,colbacktitle=warnC,
  fonttitle=\small}
\newcommand{\run}[1]{\texttt{\small #1}}
\newcommand{\ok}{\textcolor{okC}{\faCheck}}
\newcommand{\wait}{\textcolor{warnC}{\faHourglassHalf}}
\newcommand{\fig}{../../slides/figures}

\begin{document}

% ---------- Tiêu đề ----------
\begin{tcolorbox}[enhanced,colback=navy,colframe=navy,arc=3pt,left=12pt,right=12pt,top=10pt,bottom=10pt]
  \color{white}
  {\Large\bfseries Báo cáo tuần 2026-W46}\\[2pt]
  {\small 09/11 – 15/11/2026 · viết ngày 12/11/2026}\\[6pt]
  {\small\textbf{Người thực hiện:} Đỗ Trung Tuyến \quad
   \textbf{Đề tài:} Payment Intelligence Ecosystem \quad
   \textbf{Giai đoạn:} Stage G (sau lịch 6 tuần)}
\end{tcolorbox}

\begin{keypoint}
\begin{itemize}
  \item \textbf{Stage G (từ điểm số đến quyết định)} xong G01–G09: hiệu chỉnh xác suất, so sánh chính
        sách theo chi phí giả định, độ ổn định theo tuần, đóng gói tái lập, bộ chấm điểm theo lô.
        Mọi kết quả mới là \textbf{thăm dò}, chỉ dùng train + validation; tập test \textbf{không} bị mở lại.
  \item \textbf{\texttt{make demo}} chạy đầu-cuối cả hai sản phẩm trên máy chủ dự án (Docker Desktop):
        khớp số đã công bố (D-B 3/4, 0 báo nhầm/ngày; lệch chấm điểm 0).
  \item \textbf{Giao diện}: terminal demo dùng \texttt{rich} (ADR 0011); dashboard viết lại thành một trang
        có thanh điều hướng, luồng pipeline, biểu đồ.
  \item \textbf{Chất lượng}: 381 → \textbf{463 test pass} (1 skip), lint sạch; 23 commit trên 11 nhánh, chờ
        mở PR và CI.
\end{itemize}
\end{keypoint}

% =====================================================================
\section{Công việc đã làm}

\begin{small}
\renewcommand{\arraystretch}{1.25}
\begin{longtable}{@{}p{0.5cm}p{5.8cm}p{3.1cm}p{5.9cm}@{}}
\toprule
\textbf{} & \textbf{Việc} & \textbf{Nhánh / commit} & \textbf{Kiểm chứng} \\
\midrule
\endhead
\ok & G01 đặc tả protocol Stage G, chia validation thành V1 / V2 / V3; ADR 0009 & \run{g01} · \run{e2f962b} & Chủ dự án duyệt D1–D8 \\
\ok & G02 tách bảng “chất lượng xếp hạng” và “ngưỡng chốt trước” (README, báo cáo, dashboard) & \run{g02} · \run{dba37f6} & Không còn chỗ đọc 0,51 như hiệu quả triển khai \\
\ok & G03 hiệu chỉnh xác suất & \run{g03} · \run{50d01de} & \run{cal-04b3ad5944ce}; chạy 3 lần trùng từng byte \\
\ok & G04 so sánh chính sách (27 tổ hợp chi phí giả định) & \run{g04} · \run{fb53035} & \run{pol-af5df43efe54} \\
\ok & G05 độ ổn định theo tuần và nhóm & \run{g05} · \run{ed54347} & \run{stab-36887d48b356} \\
\ok & G06 báo cáo Stage G & \run{g06} · \run{82705cb} & Mọi số kèm mã lần chạy \\
\wait & G07 tách \texttt{experiment\_id} / \texttt{execution\_id} (ADR 0010), CI có PostgreSQL, kiểm sổ test-một-lần, gói \texttt{results/} & \run{g07} · \run{08c569d} & Test pass trên máy; \textbf{chờ CI xanh} \\
\ok & G08 bộ chấm điểm theo lô (\texttt{make score-model}, \texttt{make score}) & \run{g08} · \run{7ab81f6} & Lệch huấn luyện–chấm điểm 0 trên 2.000 dòng; phá hoại có chủ đích bị 2 test bắt \\
\ok & \texttt{make demo} cho hai sản phẩm; chạy được từ PowerShell; kiểm tra PostgreSQL trước khi chạy & \run{g08b}, \run{g08c} · \run{f8c1d70} & Chạy thật trên Docker Desktop, 7/7 kiểm tra đạt \\
\ok & G09 đồng bộ tài liệu; báo cáo thực tập mục 4.6 (PDF 25 trang); ADR 0002 Accepted & \run{g09} · \run{215df5f} & Đã xem từng trang mới \\
\wait & G10 đặc tả thử nghiệm độ bền (sự cố tăng dần, nạp trùng/muộn, lật nhãn) & \run{g10} · \run{c55b3ba} & \textbf{Chờ duyệt D1–D9}, chưa code \\
\ok & G11 terminal \texttt{rich} (ADR 0011); dashboard một trang; slide và báo cáo tuần LaTeX & \run{g11} · \run{4a0b68b}+ & Xem ảnh chụp sáng/tối/màn hẹp; 463 pass \\
\bottomrule
\end{longtable}
\end{small}

% =====================================================================
\section{Kết quả chính}

\begin{small}
\renewcommand{\arraystretch}{1.25}
\begin{tabularx}{\textwidth}{@{}p{3.4cm}X p{2.9cm}p{2cm}@{}}
\toprule
\textbf{Câu hỏi} & \textbf{Quan sát} & \textbf{Mã lần chạy} & \textbf{Loại} \\
\midrule
Mô hình chính (trích lại) & GBM PR-AUC test 0,408 [0,390; 0,424]; ngưỡng 0,853: precision 0,516, recall 0,371, FPR 1,24\% & \run{run-f8807cb6f780} & Test một lần (Stage E) \\
Điểm có phải xác suất? & Không: trung bình điểm 26,1\% so với 3,5\% thật ($\approx$7 lần). Platt (fit V1) đưa log loss V3 0,351 → 0,109 & \run{cal-04b3ad5944ce} & Thăm dò \\
Ngưỡng có “chuyển” được? & Ngưỡng FPR 5\% chọn ở khối trước cho FPR thực tế 3,9–5,6\% ở khối sau & \run{pol-af5df43efe54} & Thăm dò \\
Chính sách nào rẻ nhất? & Luật chi phí theo số tiền, dùng xác suất đã hiệu chỉnh: 27/27 tổ hợp chi phí \textbf{giả định} & \run{pol-af5df43efe54} & Độ nhạy, không phải ROI \\
Có ổn định theo thời gian? & PR-AUC tuần validation 0,35–0,56; PSI 4 đầu vào < 0,05 không giải thích được hai tuần cuối & \run{stab-36887d48b356} & Thăm dò \\
Chấm điểm theo lô đúng? & Điểm lô trùng điểm lúc huấn luyện (chênh 0); từ chối cột nhãn và dòng giai đoạn test & \run{score-19e94be9616d} & Kiểm thử kỹ thuật \\
Product 01 tái lập? & \texttt{make demo}: D-B 3/4 (0 báo nhầm/ngày), D-A 2/4 (1,0) trên bộ tuning, khớp công bố & \run{det-3ddc26cff253} & Tuning \\
\bottomrule
\end{tabularx}
\end{small}

\begin{figure}[H]
  \centering
  \includegraphics[width=0.92\linewidth]{\fig/p02_calibration.pdf}
  \caption{Hiệu chỉnh trên khối V3: điểm gốc thổi lên khoảng 7 lần; Platt và dịch tiên nghiệm đưa xác
  suất trung bình về gần tỷ lệ gian lận thật. PR-AUC không đổi (hiệu chỉnh đơn điệu).}
\end{figure}

\begin{figure}[H]
  \centering
  \begin{minipage}[t]{0.55\linewidth}\centering
    \includegraphics[width=\linewidth]{\fig/p02_policy.pdf}
    \caption{Tổng chi phí giả định trên V3 (mất $0{,}75\times$ số tiền, duyệt 5, nhầm 1).}
  \end{minipage}\hfill
  \begin{minipage}[t]{0.42\linewidth}\centering
    \includegraphics[width=\linewidth]{\fig/p02_stability.pdf}
    \caption{PR-AUC và tỷ lệ gian lận theo tuần (T: train, trong mẫu; V: validation).}
  \end{minipage}
\end{figure}

% =====================================================================
\section{Demo và giao diện}

\texttt{make up} rồi \texttt{make demo} chạy Product 01 (sinh dữ liệu → PostgreSQL → KPI → detector →
dashboard) và Product 02 (huấn luyện bộ chấm điểm → chấm 2.000 giao dịch). Terminal hiện sản phẩm đang
chạy, phần trăm tiến độ, bảng kết quả và danh sách kiểm tra. Stage G chạy lại bằng
\texttt{make calibration policy stability} (cần dataset); hình trong báo cáo này: \texttt{make slides-figures}.
Dashboard (\texttt{artifacts/dashboard/index.html}) là một trang HTML tĩnh mở offline:

\begin{figure}[H]
  \centering
  \includegraphics[width=0.95\linewidth]{\fig/dashboard_top.png}
  \caption{Dashboard sau \texttt{make demo}: số liệu chính và luồng pipeline của hai sản phẩm (\faCheck\ đã chạy,
  \faLock\ tập test chỉ đánh giá một lần). Phía dưới có biểu đồ theo giờ, bảng sự cố, kết quả Stage G.}
\end{figure}

% =====================================================================
\section{Chất lượng và tự kiểm tra}

\begin{minipage}[c]{0.46\linewidth}
  \includegraphics[width=\linewidth]{\fig/tests_progress.pdf}
\end{minipage}\hfill
\begin{minipage}[c]{0.5\linewidth}\small
\begin{itemize}
  \item Số test pass tăng 381 → 463 (1 skip là hook git); 7 test database chạy thật với PostgreSQL.
  \item Mỗi lệnh Stage G chạy ít nhất 2–3 lần, \texttt{results.json} trùng từng byte.
  \item Tự kiểm tra: đảo nhãn, nhìn tương lai, tính lại độc lập, phá hoại có chủ đích (ví dụ lệch
        \texttt{hour\_of\_day} 7 giờ khi chấm điểm bị 2 test bắt).
  \item Kiểm tra trên Python 3.11 bằng ruff (máy chỉ có 3.12); CI chạy 3.11 và 3.13.
\end{itemize}
\end{minipage}

\begin{warnbox}{Vi phạm quy tắc đã phát hiện và sửa trong tuần}
\small Dashboard cũ đọc \texttt{artifacts/runs/run-f8807cb6f780/test/metrics.json} để hiện số test, và
từ khi có \texttt{make demo}, mỗi lần chạy demo trên máy chủ dự án đều đọc file này. File chỉ chứa số
đã công bố (không có đánh giá lại trên tập test), nhưng việc đọc vi phạm quy tắc “không đọc thư mục
test”. \textbf{Khắc phục}: số test và số bộ report nay chỉ trích từ \texttt{results/published.json} (trong
git); thêm test cài giá trị mồi vào \texttt{runs/*/test/} để chặn tái diễn.
\end{warnbox}

% =====================================================================
\section{Hạn chế (kết quả tuần này không chứng minh)}

\begin{itemize}
  \item Không chứng minh hiệu quả khi vận hành và không phải ROI: chi phí trong G04 là giả định.
  \item “27/27” có phần đúng theo cách xây dựng: luật chi phí được chấm bằng chính hàm chi phí nó tối ưu.
  \item Validation đã bị dùng nhiều lần (lưới GBM, rò rỉ, đặc trưng lịch sử), nên số Stage G có thể lạc quan.
  \item Config Stage G được commit cùng code sau lần chạy đầu, dù viết trước và chạy lại trùng từng byte.
  \item Product 01 chỉ có 4 sự cố mỗi bộ, dữ liệu giả lập; hai sản phẩm vẫn không chung dữ liệu.
  \item CI (test database thật, job \texttt{demo}, job kiểm sổ) chưa được xác nhận xanh vì chưa mở PR.
\end{itemize}

% =====================================================================
\section{Kế hoạch tuần sau}

\begin{enumerate}
  \item Mở chuỗi PR xếp chồng (g01 → \texttt{main}, g02 → g01, …, g11 → g08c), xác nhận CI xanh, merge từ dưới lên.
  \item G10 thử nghiệm độ bền, nếu chủ dự án duyệt D1–D9.
  \item Tập trình bày bằng slide (\texttt{docs/slides/slides.pdf}) và demo 5 phút.
  \item Cập nhật hình dashboard trong báo cáo thực tập (hiện còn bản Stage F).
\end{enumerate}

% =====================================================================
\section{Cần mentor / chủ dự án quyết định}

\begin{itemize}
  \item Duyệt đặc tả G10 (\texttt{docs/architecture/stage-g10-spec.md}), đặc biệt D1 (không tăng
        \texttt{GENERATOR\_VERSION}) và D2 (nơi chạy phần cần PostgreSQL).
  \item Có cần nối hai sản phẩm thành một pipeline không? Đề xuất: không, vì không có dataset công khai
        chung khoá giao dịch; trình bày là một nền tảng, hai pipeline.
\end{itemize}


\end{document}
